\section{Einleitung}
\label{sec:einleitung}

\emph{dnchub-app} (auch bekannt als \emph{FleetOptima}) ist
ein mandantenfähiges Flottenmanagement-System, das als Open-Source-Projekt auf
GitHub verfügbar ist. Das System besteht aus einem Python/FastAPI-Backend und einem
Next.js-Frontend und ist als SaaS-Plattform konzipiert, auf der mehrere unabhängige
Organisationen ihre Fahrzeugflotten, Mitarbeiter, Werkzeuge und Verbrauchsmaterialien
verwalten können.

\subsection{Aufgabenstellung}

Im Rahmen des Moduls CySL an der FHNW wurde ein nicht-trivialer,
sicherheitsrelevanter Patch für ein FOSS-Projekt auf GitHub als
Prüfungsleistung gewählt. Die Wahl fiel auf eine kritische IDOR-Schwachstelle
(GitHub Issue \#5, CVSS 9.3) \cite{github_issue_5}, die bei einem Code-Review
des Repository-Owners identifiziert wurde und direkt mit einer zweiten
Schwachstelle (GitHub Issue \#14) \cite{github_issue_14} zusammenhängt.

Beide Issues sind miteinander verknüpft und können, laut Empfehlung des
Repository-Owners, in einem einzigen Pull Request behoben werden. Der Dozent
schätzt den Aufwand auf circa drei Tage (mit sauberem Testing eher vier Tage),
was einer realistischen Einschätzung entspricht.

\subsection{Ziele}

Das Ziel war kein punktueller Fix, der eine Fehlermeldung zum Verstummen
bringt. Gesucht war eine Lösung an der Wurzel des Problems, deren Wirksamkeit
sich mit Tests beweisen lässt. Im Einzelnen:

\begin{itemize}
      \item Analyse der verwundbaren Codebasis von \texttt{BaseService}
            bis zu allen abhängigen Aufrufern
      \item Sicherer Fix für \texttt{get()} und
            \texttt{get\_or\_404()} mit \texttt{organization\_id}-Scoping
            und \texttt{deleted\_at}-Filterung; danach alle betroffenen
            Aufrufer im Backend anpassen
      \item Automatisierte Sicherheitstests, die die Schwachstelle auf
            dem Originalstand reproduzieren und den Fix verifizieren
      \item Vollständige Dokumentation des Vorgehens
\end{itemize}

\subsection{Aufbau des Berichts}

Kapitel \ref{sec:grundlagen} legt die theoretischen Grundlagen zu IDOR,
Multi-Tenant-Architektur und Soft-Delete-Mustern.
Kapitel \ref{sec:analyse} zeigt, wo die Schwachstellen im Code stecken und
wie viel des Backends betroffen ist (Phase 1).
Kapitel \ref{sec:implementierung} beschreibt den Patch und seine Ausbreitung
auf alle Aufrufer (Phase 2).
Kapitel \ref{sec:testing} dokumentiert Teststrategie und Ergebnisse (Phase 3).
Kapitel \ref{sec:schluss} zieht Bilanz.
